\documentclass[12pt,a4paper]{article}
\usepackage[T1]{fontenc}
\usepackage{lmodern}
\usepackage{amsmath,amssymb,mathtools}
\usepackage{booktabs,longtable}
\usepackage[margin=25mm]{geometry}
\usepackage{microtype}
\usepackage[hidelinks]{hyperref}
\setlength{\emergencystretch}{2em}
\newcommand{\bp}{\mathrm{bp}}
\newcommand{\BuildResult}{Successfully compiled with latexmk 4.88 and MiKTeX-pdfTeX 4.27 (MiKTeX 26.5).}
\title{Independent Mathematical Verification\\
of the Fixed-Income Analytics Engine}
\author{Technical verification record}
\date{14 September 2026}
\begin{document}
\maketitle
\begin{abstract}
The implemented fixed-rate bond, Nelson--Siegel--Svensson (NSS), and
continuous-zero curve-risk mathematics were reviewed independently against
their stated assumptions. Analytical derivations, source inspection, existing
tests, and additional numerical calculations agree. No mathematical defect
was found within the reviewed conventions. The baseline suite passed all
234 tests; all four illustrative command-line examples completed.
Duration and convexity checks exhibit the expected finite-difference error
behaviour. Independently generated noiseless NSS observations fit to near
floating-point precision. Key-rate weights form a partition of unity;
the small difference between summed key-rate and parallel full-repricing
DV01 is explained by a cubic term.
This is a verification report for later academic use, not a dissertation
or an unrestricted certification of all possible floating-point inputs.
\end{abstract}

\section{Review basis, scope, and reproducibility}
The numerical review was executed on 13 September 2026; this report was
finalized on 14 September 2026. The audited source revision is:
\begin{center}\small
\texttt{7eee019c2b186b95791bba42a7dc8edf5c376047}
\end{center}
The initial branch was \texttt{main}, with an empty
\texttt{git status --short}. The remote was
\url{https://github.com/REMEDIOSSAMUEL/fixed-income-analytics-engine.git}.
The six inspected commits, newest first, were \texttt{7eee019},
\texttt{33ce60a}, \texttt{956337d}, \texttt{b8ed944},
\texttt{677efea}, and \texttt{ef25a20}.

Before mathematical work, the complete contents of \texttt{AGENTS.md},
\texttt{README.md}, \texttt{pyproject.toml}, all five tracked Python files
under \texttt{src/fixed\_income/}, and all three files under \texttt{tests/}
were read. The illustrative CSV and existing ignore rules were also inspected.
Implementation, documentation, and tests were treated as separate evidence:
a test that repeats the implementation cannot by itself establish the
underlying financial identity.

No financial source, tests, public interface, dependency declaration, or
CLI command was changed. Additional calculations used stdin-fed Python;
the full scripts are reproduced in Appendix~\ref{sec:scripts}, without
creating a permanent Python verification module. Standard-library
\texttt{Decimal} supplied an independent 90-digit NSS reference.
Reported numerical values come from actual execution, not the README.

\subsection{Assumptions and the scope of exactness}
A bond has fixed, nonnegative coupons, positive redemption, regular annual
or semiannual periods, and maturity strictly after settlement.
There are no irregular coupons, issue-date constraints, business-day
adjustments, ex-coupon rules, credit events, optionality, or changing cash
flows. Schedules are anchored to maturity; nonexistent calendar days clip
to the end of the corresponding month. There is no additional end-of-month
rule. The settlement-date coupon is outside the future payment set.

Nominal annual YTM discounting uses a strictly positive periodic
accumulation factor. Curve-risk discounting separately uses explicitly
interpreted continuously compounded zero rates and actual days divided by
365. Negative yields and zero rates are compatible with the mathematics
when the expressions are valid. Shocks hold settlement and cash flows fixed.
Curve and shock callables are deterministic.

An \emph{exact identity} below means equality over real arithmetic under these
assumptions. Numerical powers, exponentials, sums, and optimizations introduce
rounding or approximation. Tests and grids provide evidence at their inputs;
analytical arguments establish general identities. This review does not
establish universal market day-count applicability, global calibration
optimality, parameter uniqueness, arbitrage freedom, or accuracy for every
extreme representable input.

\subsection{Execution environment}
Execution used native Windows PowerShell and the repository virtual environment.
Python was 3.12.10, 64-bit, on Windows 11
(\texttt{10.0.26200}). Installed versions were NumPy 2.5.3,
SciPy 1.18.1, pandas 3.0.5, matplotlib 3.11.2,
python-dateutil 2.9.0.post0, and pytest 9.1.1.
The package's \texttt{direct\_url.json} reported
\texttt{"editable": true}; import resolved to this repository's
\texttt{src/fixed\_income/\_\_init\_\_.py}.
The existing editable installation and package import were thus verified;
the environment was not reinstalled. Dependencies are not locked by the
project, so another numerical environment may round or optimize differently.

\section{Notation and units}
\small
\begin{longtable}{@{}p{0.20\textwidth}p{0.74\textwidth}@{}}
\toprule
Symbol & Definition and convention\\
\midrule
\endhead
$F$ & Face/redemption amount in currency units; default 100.\\
$c$ & Annual coupon rate as a decimal.\\
$m$ & Coupon payments per year, 1 or 2; $1/m$ converts coupon periods to years.\\
$C$ & Currency coupon cash flow per payment period.\\
$y$ & Nominal annual decimal YTM compounded $m$ times per year.\\
$\alpha,w$ & Accrued and remaining fractions of the current coupon period; $w=1-\alpha$.\\
$N,i$ & Number of strictly future payments and payment index $1\leq i\leq N$.\\
$q_i$ & Fractional coupon-period count from settlement to payment $i$.\\
$T_i=q_i/m$ & Payment time in coupon-based years for YTM analytics.\\
$CF_i,PV_i$ & Currency cash flow and its YTM-discounted present value.\\
$P_{\rm dirty},P_{\rm clean},AI$ & Dirty price, clean price, accrued interest, in currency for the supplied face value.\\
$P(y),P',P''$ & Dirty YTM price and derivatives with respect to annual decimal $y$.\\
$D_{\rm Mac},D_{\rm Mod},\mathcal C$ & Macaulay and modified durations (years), and normalized convexity (conventionally years squared).\\
$h$ & Decimal-rate perturbation; DV01 uses $h=10^{-4}$.\\
$T$ & Generic curve maturity in years.\\
$r_{\rm NSS}(T)$ & NSS decimal rate; its compounding convention is supplied by the caller.\\
$\beta_0,\ldots,\beta_3$ & NSS coefficients in decimal rate units.\\
$\tau_1,\tau_2$ & Strictly positive NSS decay scales in years.\\
$x,L_1,L_2$ & Dimensionless ratio $x=T/\tau$ and NSS factor loadings.\\
$\theta$ & NSS vector $(\beta_0,\beta_1,\beta_2,\beta_3,\tau_1,\tau_2)$.\\
$n,e_i$ & Observation count and observed-minus-fitted decimal-rate residual.\\
$t_i$ & Actual settlement-to-payment calendar days divided by 365, in curve-risk years.\\
$r(t_i),DF(t_i),V_i$ & Continuous annual decimal zero rate, discount factor, and currency zero-discounted cash-flow PV.\\
$s,P(s)$ & Parallel decimal zero-rate shift and corresponding dirty price.\\
$k_j,b_j(T),J$ & Ordered key maturities, dimensionless key basis, and number of keys.\\
$P_j^\pm$ & Fully repriced dirty price under $\pm h b_j$ zero-rate changes.\\
$\Delta$ & Sum of key-rate DV01s minus parallel DV01, in price units.\\
$\epsilon$ & Binary64 machine precision, approximately $2.22\,10^{-16}$.\\
$O(h^p)$ & Remainder bounded locally by a constant times $|h|^p$, with other inputs fixed.\\
\bottomrule
\end{longtable}
\normalsize
The two times $T_i$ and $t_i$ distinguish the YTM and zero-curve conventions.

\section{Fixed-rate bond foundations and fractional-period pricing}
\subsection{Coupons, accrual, and clean price}
The annual contractual coupon amount is $Fc$. Dividing it equally among
$m$ regular payments gives the exact identities
\begin{equation}
 C=\frac{Fc}{m},\qquad CF_i=C+F\,\mathbf 1_{\{i=N\}},
\end{equation}
where the indicator $\mathbf 1_{\{i=N\}}$ is one at redemption and zero
otherwise. Currency times annual coupon rate divided by payments per year
is currency per payment.

Let $d_-$ be the coupon date on or immediately before settlement $d_s$,
and $d_+$ the first strictly later coupon date. The project defines
\begin{equation}
 \alpha=\frac{\operatorname{days}(d_s-d_-)}
                   {\operatorname{days}(d_+-d_-)},\quad
 0\leq\alpha<1,\qquad AI=C\alpha.
\end{equation}
The calendar-day ratio is dimensionless. Linear accrual over this particular
coupon period is an assumption, not a consequence of discounting.
This Actual/Actual-within-period convention does not implement every market's
day-count or accrued-interest rule. By the project's clean-price definition,
\begin{equation}
 P_{\rm clean}=P_{\rm dirty}-AI,\qquad
 P_{\rm dirty}=P_{\rm clean}+AI.
\end{equation}
These are exact accounting identities under the chosen accrual rule.
The amounts are not implicitly normalized to 100.

The source implements accrual in \texttt{bonds.py:123--130};
\texttt{clean\_price} at line 168 subtracts it. Existing tests count 90
elapsed days in a 181-day period and check an exactly half-accrued leap-year
period.

\subsection{Derivation of between-coupon pricing}
Define the periodic accumulation factor $B=1+y/m>0$.
For a constant YTM under the fractional-period convention, accumulation
over $q$ coupon periods multiplies an amount by $B^q=\exp(q\log B)$;
discounting is its inverse. This extension of periodic compounding to
fractional periods is an explicit assumption.
Regular future payments are one coupon period apart, so
\begin{equation}
 w=1-\alpha,\qquad q_1=w,\qquad q_i=w+i-1,\qquad T_i=\frac{q_i}{m}.
\end{equation}
Linearity of valuation of fixed cash flows then yields
\begin{equation}\label{eq:ytm}
 P(y)=P_{\rm dirty}(y)=\sum_{i=1}^{N}CF_i B^{-q_i},
 \qquad PV_i=CF_i B^{-q_i}.
\end{equation}
Exponents are dimensionless period counts and need not be integers:
settlement can occur partway through a period. $T_i$ is not generally
actual days divided by 365.

In \texttt{bonds.py:139--159}, the zero-based loop index is $i-1$:
\texttt{q = w + index}. The internal schedule has one preceding/on-settlement
date before the future dates. Therefore \texttt{len(dates)-1} counts future
payments and \texttt{index == len(dates)-2} identifies the last one for
principal addition. This agrees with (\ref{eq:ytm}) without an off-by-one error.

\subsection{Settlement exactly on a coupon date}
The backward schedule at \texttt{bonds.py:94--107} stops at a date less
than or equal to settlement. The public schedule removes that first date.
On a coupon date, $d_-=d_s$, $\alpha=0$, $AI=0$, $w=1$, and $q_1=1$.
That day's coupon is excluded, and the next payment is one full period away.

Actual checks with maturity 15 January 2027, $F=100$, $c=0.04$, $m=2$:
\begin{center}\small
\begin{tabular}{@{}llllrr@{}}
\toprule
Settlement & Previous & Next & $N$ & $AI$ & $q_1$\\
\midrule
2025-07-14 & 2025-01-15 & 2025-07-15 & 4 & 1.988950276243 & 0.005524861878\\
2025-07-15 & 2025-07-15 & 2026-01-15 & 3 & 0 & 1\\
2025-07-16 & 2025-07-15 & 2026-01-15 & 3 & 0.010869565217 & 0.994565217391\\
\bottomrule
\end{tabular}
\end{center}
The future-payment count decreases precisely on the coupon date.
Month-end anchoring and annual leap-day tests were also inspected: each
date is calculated from maturity, not recursively from an already clipped date.

\section{Par-bond identity: geometric-series proof}
On a coupon date $q_i=i$. Set $c=y$, $a=y/m$, $v=(1+a)^{-1}$, so $C=Fa$.
Then
\[
 P=Fa\sum_{i=1}^{N}v^i+Fv^N.
\]
For $a\ne0$, multiplying the geometric sum by $1-v$ gives
\[
 (1-v)\sum_{i=1}^{N}v^i=v-v^{N+1}=v(1-v^N).
\]
Since $1-v=av$,
\begin{equation}
 \sum_{i=1}^{N}v^i=\frac{1-v^N}{a},\qquad
 P=F(1-v^N)+Fv^N=F.
\end{equation}
For $a=0$, coupons are zero and undiscounted redemption alone gives $P=F$.
These identities are exact for a regular bond on a coupon date with matching
coupon and nominal YTM compounding. They do not assert dirty par between coupons.

Independent sums used settlement 15 January 2025, maturity 15 January 2030,
$F=100$ and $c=y=0.04$. Annual ($N=5$) and semiannual ($N=10$) cases each
returned \texttt{99.99999999999999}; the engine returned the same values.
The approximately $1.4\,10^{-14}$ distance from 100 is rounding. Accrual is
zero, so clean and dirty par agree.

\section{First derivative and duration}
Hold cash flows, dates, $q_i$, and $m$ fixed. For $B>0$ the real-power
chain rule applies even to fractional $q_i$:
\[
 \frac{d}{dy}B^{-q_i}=-q_i B^{-q_i-1}\frac1m.
\]
Termwise differentiation of the finite sum gives the exact identity
\begin{equation}\label{eq:first}
 P'(y)=-\sum_i CF_i\frac{q_i}{m}B^{-q_i-1}
       =-\frac1B\sum_i T_iPV_i.
\end{equation}
Macaulay duration is the PV-weighted payment time,
\[
 D_{\rm Mac}=\frac{\sum_i T_iPV_i}{P}.
\]
Normalizing (\ref{eq:first}) derives modified duration:
\begin{equation}\label{eq:mod}
 D_{\rm Mod}=-\frac{P'}P=\frac{D_{\rm Mac}}{1+y/m}.
\end{equation}
No integer-exponent assumption was used. $D_{\rm Mac}$ is in years, and
division by the dimensionless $B$ preserves that time convention.
Modified duration multiplies an annual decimal yield change, not a
percentage-number change.

The normalized weighted sum at \texttt{bonds.py:175--182} implements
$D_{\rm Mac}$; lines 184--190 divide it by $B$.
At fixed settlement clean and dirty prices have the same unnormalized
YTM derivatives, since $AI$ is yield-independent. Normalizing by clean
price would nevertheless define a different duration from this API.

\section{Independent numerical duration and convexity checks}
\subsection{Independent reference cash flows}
The numerical bond has settlement 15 April 2025, maturity 15 January 2030,
$F=100$, $c=0.04$, $m=2$, $y=0.045$.
The independent script explicitly constructed ten dates: 15 July 2025,
15 January and 15 July in 2026--2029, and 15 January 2030.
It used nine cash flows of 2 and a final cash flow of 102, without calling
engine pricing or analytics methods to produce the reference price function.

Independent day counts give $\alpha=90/181$, with values
\begin{align*}
 \alpha&=0.4972375690607735,& w&=0.5027624309392265,\\
 q_N&=9.502762430939226,& AI&=0.994475138121547.
\end{align*}
\begin{center}\small
\begin{tabular}{@{}lrr@{}}
\toprule
Quantity & Independent calculation & Engine\\
\midrule
Dirty price & 98.87131308203281 & 98.87131308203281\\
Clean price & 97.87683794391127 & 97.87683794391127\\
$D_{\rm Mac}$ & 4.326726840097139 & 4.326726840097138\\
$D_{\rm Mod}$ & 4.231517692026542 & 4.231517692026541\\
$\mathcal C$ & 21.110444849612590 & 21.110444849612588\\
\bottomrule
\end{tabular}
\end{center}
Independent analytical derivatives were $P'=-418.3757105405171$ and
$P''=2087.2174020270336$, per unit annual decimal YTM and its square.

\subsection{Duration finite differences and their order}
Taylor expansion gives
\begin{align}
 P(y+h)-P(y-h)&=2hP'(y)+\frac{h^3}{3}P'''(y)+O(h^5),\\
 -\frac{P(y+h)-P(y-h)}{2hP(y)}
 &=D_{\rm Mod}-\frac{h^2P'''(y)}{6P(y)}+O(h^4).
\end{align}
Here $P'''$ denotes the third YTM derivative. This central estimate has
second-order truncation error $O(h^2)$; finite-$h$ equality is an approximation.
The denominator is the independently calculated unshifted dirty price.

\begin{center}\small
\begin{tabular}{@{}rrr@{}}
\toprule
$h$ (decimal YTM) & Central duration (years) & Signed error (years)\\
\midrule
$10^{-2}$ & 4.23346867529047 & $+1.950983264\,10^{-3}$\\
$10^{-3}$ & 4.23153719801642 & $+1.950598988\,10^{-5}$\\
$10^{-4}$ & 4.23151788708521 & $+1.950586661\,10^{-7}$\\
$10^{-5}$ & 4.23151769399938 & $+1.972840558\,10^{-9}$\\
$10^{-6}$ & 4.23151769257645 & $+5.499058986\,10^{-10}$\\
$10^{-7}$ & 4.23151769415749 & $+2.130945198\,10^{-9}$\\
$10^{-8}$ & 4.23151766397402 & $-2.805252652\,10^{-8}$\\
\bottomrule
\end{tabular}
\end{center}
Error means estimate minus independent analytical value. A tenfold step
reduction initially reduces error by about 100. At $h=10^{-5}$ relative
error is about $4.7\,10^{-10}$. Subtracting nearby prices eventually
amplifies roundoff, with a normalized first-difference error scale of
order $\epsilon/h$ up to problem-dependent constants. An intermediate
step is preferable to an arbitrarily tiny one.

\subsection{Convexity: independent second differentiation}
Differentiating (\ref{eq:first}) again contributes $(-q_i-1)/m$ and
reverses the sign:
\begin{equation}\label{eq:second}
 P''(y)=\sum_i CF_i\frac{q_i(q_i+1)}{m^2}B^{-q_i-2}.
\end{equation}
This identity is exact for real fractional $q_i$ and $B>0$. The project defines
\begin{equation}
 \mathcal C=\frac{P''}P
 =\sum_i\frac{PV_i}P\frac{q_i}{mB}\frac{q_i+1}{mB}.
\end{equation}
Convexity is the normalized second derivative with respect to annual
decimal YTM, conventionally in years squared. Both $1/m$ factors convert
period counts to the corresponding time units. No factor $1/2$ enters
$\mathcal C$ itself; it belongs in the price expansion for small decimal
yield change $\delta y$:
\begin{equation}
 \frac{P(y+\delta y)-P(y)}{P(y)}
 =-D_{\rm Mod}\delta y+\frac12\mathcal C(\delta y)^2+O((\delta y)^3).
\end{equation}
The weighted product at \texttt{bonds.py:192--206} is algebraically
(\ref{eq:second}). Positive cash flows and positive $q_i$ imply
positive convexity.

Writing $P^{(4)}$ for the fourth derivative, central second differences obey
\begin{equation}
 \frac{P(y+h)-2P(y)+P(y-h)}{h^2}
 =P''(y)+\frac{h^2}{12}P^{(4)}(y)+O(h^4).
\end{equation}
The derivative and its normalized value have $O(h^2)$ truncation error.
\begin{center}\small
\begin{tabular}{@{}rrrr@{}}
\toprule
$h$ & Central second derivative & Normalized value & Error in $\mathcal C$\\
\midrule
$10^{-2}$ & 2087.80194712077 & 21.1163570305629 & $+5.912180950\,10^{-3}$\\
$10^{-3}$ & 2087.22324657629 & 21.1105039623023 & $+5.911268972\,10^{-5}$\\
$10^{-4}$ & 2087.21744030527 & 21.1104452367647 & $+3.871520668\,10^{-7}$\\
$10^{-5}$ & 2087.21758099273 & 21.1104466596998 & $+1.810087163\,10^{-6}$\\
$10^{-6}$ & 2087.20507544058 & 21.1103201765798 & $-1.246730328\,10^{-4}$\\
$10^{-7}$ & 2069.10044653341 & 20.9272071143295 & $-1.832377353\,10^{-1}$\\
$10^{-8}$ & 2131.62820728030 & 21.5596227139383 & $+4.491778643\,10^{-1}$\\
\bottomrule
\end{tabular}
\end{center}
Convergence at larger steps precedes cancellation-driven deterioration.
At $10^{-4}$ the relative error is about $1.8\,10^{-8}$.
Three nearly cancelling price terms are divided by $h^2$, amplifying
roundoff on a scale of order $\epsilon/h^2$.
Balancing schematic $h^2$ truncation against rounding suggests
problem-dependent scales near $\epsilon^{1/3}$ for central first differences
and $\epsilon^{1/4}$ for central second differences. These are scaling
estimates, not universal optimal-step formulas.

\section{YTM DV01: repricing versus duration approximation}
One basis point is exactly
\[
 1\bp=0.01\text{ percentage point}=0.0001=10^{-4}
\]
in decimal rate units. The positive project DV01 convention is a currency
price response for one basis point, defined by exact central full repricing:
\begin{equation}
 {\rm DV01}_{\rm YTM}=\frac{P(y-h)-P(y+h)}2,\qquad h=10^{-4}.
\end{equation}
It is not divided by $h$ and is not a derivative per unit decimal rate.
Taylor expansion derives
\begin{equation}
 \frac{P(y-h)-P(y+h)}2
 =-hP'(y)-\frac{h^3}{6}P'''(y)+O(h^5)
 =hPD_{\rm Mod}+O(h^3).
\end{equation}
The first-order approximation has $O(h^3)$ absolute error and generally
$O(h^2)$ relative error if $P'\ne0$. Positive future cash flows imply
$P'<0$: lowering yield raises price, so down-minus-up DV01 is positive.
Also $P'''<0$, giving a positive leading cubic correction.

The implementation at \texttt{bonds.py:208--217} fully reprices both
shifted yields, each of which must satisfy yield and numerical-price limits.
For the independently constructed bond:
\begin{center}\small
\begin{tabular}{@{}lr@{}}
\toprule
Quantity (currency for $F=100$) & Actual value\\
\midrule
Independent $P(y-h)$ & 98.91316109110264\\
Independent $P(y+h)$ & 98.82948594513739\\
Implemented DV01 & 0.04183757298262236\\
Independent central repricing & 0.04183757298262236\\
Independent $PD_{\rm Mod}10^{-4}$ & 0.04183757105405172\\
Central minus duration approximation & $1.928570643461303\,10^{-9}$\\
\bottomrule
\end{tabular}
\end{center}
Agreement of full repricing is exact to the displayed precision.
The difference from the duration approximation has the predicted sign
and is a finite-bump effect, not a conversion defect.

\section{NSS model, analytical limits, and numerical stability}
\subsection{Model and independent loading expansions}
For $T>0$ and $\tau>0$ define
\begin{equation}
 L_1(T;\tau)=\frac{1-\exp(-T/\tau)}{T/\tau},\qquad
 L_2(T;\tau)=L_1(T;\tau)-\exp(-T/\tau).
\end{equation}
These are dimensionless because $T/\tau$ is dimensionless. The model is
\begin{equation}\label{eq:nss}
 \begin{split}
 r_{\rm NSS}(T)={}&\beta_0+\beta_1 L_1(T;\tau_1)
                  +\beta_2 L_2(T;\tau_1)\\
                 &+\beta_3 L_2(T;\tau_2).
 \end{split}
\end{equation}
This is a parametric specification, not a bond-pricing identity.
The four columns in \texttt{curves.py:45--50} are exactly
$1,L_1(T;\tau_1),L_2(T;\tau_1),L_2(T;\tau_2)$;
\texttt{yield\_rate} at line 77 multiplies them by the betas.
Betas are decimal rates and taus are years. The API does not assign a
compounding convention, bootstrap, or convert par yields into zeros.

Set $x=T/\tau$. Subtracting the exponential series from 1, dividing by
$x$, and then subtracting the exponential again derives
\begin{align}
 e^{-x}&=1-x+\frac{x^2}{2}-\frac{x^3}{6}
               +\frac{x^4}{24}-\frac{x^5}{120}+\cdots,\\
 L_1&=1-\frac{x}{2}+\frac{x^2}{6}-\frac{x^3}{24}
              +\frac{x^4}{120}-\frac{x^5}{720}+O(x^6),\\
 L_2&=\frac{x}{2}-\frac{x^2}{3}+\frac{x^3}{8}
              -\frac{x^4}{30}+\frac{x^5}{144}+O(x^6).
\end{align}
The coefficient of $x^n$ in $L_2$, $n\geq1$, is
$(-1)^{n+1}n/(n+1)!$, independently confirming the polynomial coefficients.
The constant terms establish the exact limits
\begin{equation}
 L_1(0^+;\tau)=1,\quad L_2(0^+;\tau)=0,\quad
 r_{\rm NSS}(0^+)=\beta_0+\beta_1.
\end{equation}
These are limits, not evaluation at zero; the API rejects $T=0$.
For $x\to\infty$, $e^{-x}\to0$ and $0<L_1<1/x\to0$.
Consequently $L_2\to0$, establishing the exact long limit
\begin{equation}
 r_{\rm NSS}(\infty)=\beta_0.
\end{equation}
For fixed positive taus the asymptotic expression is
\begin{equation}
 r_{\rm NSS}(T)=\beta_0+
 \frac{(\beta_1+\beta_2)\tau_1+\beta_3\tau_2}{T}
 +\text{exponentially decaying terms}.
\end{equation}

Thus $\beta_0$ is the asymptotic level and $\beta_1$ sets the short-limit
offset. Calling $\beta_1$ a slope coefficient does not mean it equals
the curve's derivative at zero: the expansions instead give
\[
 r_{\rm NSS}'(0^+)
 =\frac{-\beta_1+\beta_2}{2\tau_1}+\frac{\beta_3}{2\tau_2}.
\]
For $x>0$, $e^x>1+x$ implies $L_2=[1-(1+x)e^{-x}]/x>0$.
Since it vanishes at both ends, $L_2$ has an interior maximum.
Its beta coefficient contributes a hump or trough according to sign,
but overlapping factors need not create separate visible humps.
Taus rescale loading locations and widths; they do not assert market
turning points. No monotonicity or arbitrage constraints follow from NSS alone.

\subsection{Cancellation and the stable implementation}
For small $x$, 1 and $e^{-x}$ are near 1, while their difference is of
order $x$. Subtracting rounded near-equal numbers can give relative error
of order $\epsilon/x$. Eventually $e^{-x}$ rounds to 1 and the naive
numerator becomes zero. Even stable $L_1$ followed by subtraction of
$e^{-x}$ can lose the small $L_2\sim x/2$.

At \texttt{curves.py:29--42}, $x<10^{-3}$ selects separate Taylor
polynomials through $x^5$ for both loadings, with $O(x^6)$ absolute
truncation error. The next coefficients are $1/5040$ and $-1/840$;
near the threshold their term magnitudes are about $2\,10^{-22}$
and $1.2\,10^{-21}$, below ordinary rounding at the loading scales.
These are error estimates, not exact equality of polynomial and loading.
Otherwise $L_1=-\operatorname{expm1}(-x)/x$ and $L_2=L_1-e^{-x}$.
The computational switch can produce rounding-level discrepancies;
the underlying mathematical functions are continuous.
Underflow of positive $T/\tau$ uses limiting polynomial values;
an overflowing ratio uses limiting zero loadings.

Independent checks used
$\theta=(0.04,-0.025,0.03,-0.012,1.4,6.5)$, whose short limit is 0.015.
The reference used the original exponential expression at 90-digit
decimal precision, then binary64 conversion. All outputs were finite:
\begin{center}\small
\begin{tabular}{@{}rrr@{}}
\toprule
$T$ (years) & Engine decimal rate & Engine minus reference\\
\midrule
$10^{-24}$ & 0.014999999999999999 & 0\\
$10^{-18}$ & 0.014999999999999999 & 0\\
$10^{-12}$ & 0.015000000000018719 & $-1.735\,10^{-18}$\\
$10^{-8}$ & 0.015000000187197802 & 0\\
0.000999 & 0.015018693943217502 & 0\\
0.001 & 0.015018712648743639 & 0\\
0.001001 & 0.015018731354255518 & 0\\
\bottomrule
\end{tabular}
\end{center}
Zero means equal binary64 values, not zero error against exact arithmetic.
Isolating $L_2(T;1)$ avoids a nonzero curve level hiding loading error:
\begin{center}\small
\begin{tabular}{@{}rrr@{}}
\toprule
$T$ & Engine $L_2(T;1)$ & Relative reference error\\
\midrule
$10^{-24}$ & $4.9999999999999996\,10^{-25}$ & 0\\
$10^{-18}$ & $5.0000000000000004\,10^{-19}$ & 0\\
$10^{-12}$ & $4.9999999999966667\,10^{-13}$ & 0\\
$10^{-8}$ & $4.9999999666666670\,10^{-9}$ & 0\\
0.000999 & 0.00049916745759218161 & 0\\
0.001 & 0.00049966679163315764 & $-3.656\,10^{-13}$\\
0.001001 & 0.00050016612500880520 & $+4.465\,10^{-13}$\\
\bottomrule
\end{tabular}
\end{center}
Isolating all three nonconstant columns on the ten-point small-maturity
set in Appendix~\ref{sec:scripts}, including both taus' switch regions,
gave maximum relative errors $2.222669457375284\,10^{-16}$ for $L_1(T;1)$
and $4.465429311674831\,10^{-13}$ for each $L_2$ column.
Naive binary64 $(1-\exp(-T))/T$, by contrast, gave zero at $10^{-24}$
and $10^{-18}$, and 0.9999778782798785 at $10^{-12}$.

Long-maturity values approach $\beta_0=0.04$:
\begin{center}\small
\begin{tabular}{@{}rrr@{}}
\toprule
$T$ (years) & Engine rate & Rate minus $\beta_0$\\
\midrule
100 & 0.03929000266119859 & $-7.099973388014075\,10^{-4}$\\
$10^4$ & 0.039992900000000005 & $-7.099999999995998\,10^{-6}$\\
$10^8$ & 0.03999999929 & $-7.100000032345122\,10^{-10}$\\
$10^{12}$ & 0.039999999999929 & $-7.099876242477876\,10^{-14}$\\
\bottomrule
\end{tabular}
\end{center}
The asymptotic numerator is
$(\beta_1+\beta_2)\tau_1+\beta_3\tau_2=-0.071$, explaining
the approximately $-0.071/T$ deviations.

\section{NSS calibration and independent synthetic fit}
For observed pairs $(T_i,r_i^{\rm obs})$, define
\begin{equation}
 e_i(\theta)=r_i^{\rm obs}-r_{\rm NSS}(T_i;\theta),\qquad
 \min_\theta\sum_{i=1}^{n}e_i(\theta)^2.
\end{equation}
The implemented search has $\tau_1,\tau_2\in[0.01,100]$ years and
unbounded betas. Observations have equal weight, including repeated tenors.
At least six distinct positive maturities are required. The exact definitions are
\begin{equation}
 {\rm RMSE}=\sqrt{\frac1n\sum_i e_i^2},\qquad
 {\rm RMSE}_{\bp}=10{,}000\,{\rm RMSE}.
\end{equation}
Residuals and RMSE are decimal rates; the objective has squared decimal-rate
units. Scaling optimization residuals by 10,000 multiplies the squared
objective by $10^8$ without changing its mathematical minimizers.
SciPy's linear-loss cost has an additional constant $1/2$, also immaterial
to the minimizers.

At fixed taus, linear least squares initializes betas. Joint fitting is
nonlinear because taus enter exponential arguments.
At \texttt{curves.py:139--214}, the six deterministic starting pairs are
$(0.5,2)$, $(1,5)$, $(2,10)$, $(5,1)$, $(10,3)$, $(3,15)$ years.
The optimizer uses Jacobian-based scaling, tolerances $10^{-10}$, and
\texttt{max\_nfev=3000} per run.
Only successful, finite, bounded candidates with finite squared residual
sums are eligible. The lowest unscaled sum wins; exact ties retain the
earlier start. Residual signs and reported RMSE agree with the definitions.

The installed SciPy signature confirms the default Jacobian
\texttt{jac='2-point'}. A one-sided difference
$[e(\theta+\delta u)-e(\theta)]/\delta$, for parameter-coordinate unit vector
$u$, approximates that Jacobian column with $O(\delta)$ truncation error
for smooth residuals, plus cancellation error. This is separate from
analytical bond derivatives. Reported \texttt{nfev} describes the selected
run and excludes extra numerical-Jacobian residual calls; it is not the
total work over starts.

Multiple starts reduce dependence on one local solution but do not prove
global optimality. Excellent curve agreement does not imply unique
six-parameter identification. If $\beta_1=\beta_2=\beta_3=0$, both taus
are unidentifiable from the constant curve. If $\tau_1=\tau_2$, only
$\beta_2+\beta_3$ multiplies their shared loading. Nearby degeneracies
and limited tenor coverage can cause weak identification.

\subsection{Independent noiseless observations}
The 40 maturities were
$T_i=0.05(40/0.05)^{(i-1)/39}$ years, $i=1,\ldots,40$, generated by
\texttt{np.geomspace}. The known vector was
$(0.04,-0.025,0.03,-0.012,1.4,6.5)$.
Observed rates came from the independent 90-digit exponential formula,
then binary64 conversion; the engine did not generate them.
Noiseless means no added observation noise, with representation rounding
still present. Public \texttt{fit\_nss} fitted these observations.
A further 1,001-point geometric grid in the same range tested between-node
curve agreement.
\begin{center}\small
\begin{tabular}{@{}lr@{}}
\toprule
Metric & Actual result\\
\midrule
Maximum absolute observation rate error & $6.938893903907228\,10^{-18}$\\
Independently recomputed RMSE & $3.291406376790657\,10^{-18}$\\
API RMSE & $3.291406376790657\,10^{-18}$\\
RMSE in basis points & $3.291406376790657\,10^{-14}$\\
Maximum grid rate error (1,001 points) & $1.387778780781446\,10^{-17}$\\
Maximum residual-vector discrepancy & 0\\
Successful starts / attempted & 6 / 6\\
Selected optimizer status / evaluations & 3 / 9\\
\bottomrule
\end{tabular}
\end{center}
The fitted parameters were:
\begin{center}\small
\begin{tabular}{@{}lr@{}}
\toprule
Parameter & Fitted value\\
\midrule
$\beta_0$ & 0.039999999999999966\\
$\beta_1$ & -0.024999999999999967\\
$\beta_2$ & 0.030000000000000134\\
$\beta_3$ & -0.011999999999999962\\
$\tau_1$ (years) & 1.4000000000000024\\
$\tau_2$ (years) & 6.499999999999933\\
\bottomrule
\end{tabular}
\end{center}
Parameter proximity happened in this run; exact recovery was not required.
This verifies sampled curve agreement, not a global optimum or
extrapolation guarantee.

\section{Continuous-zero pricing and parallel shock monotonicity}
\subsection{Discount factors and the separate time convention}
Let $r(t_i)$ be the continuous annual decimal zero rate for payment time
$t_i$. Continuous compounding gives the maturity-specific accumulation
factor $\exp(r(t_i)t_i)$; inversion yields the exact discount factor
\begin{equation}
 DF(t_i)=e^{-r(t_i)t_i},\qquad
 P_{\rm zero}=\sum_i CF_i e^{-r(t_i)t_i}=\sum_i V_i.
\end{equation}
Here $V_i=CF_iDF(t_i)$ is the zero-discounted currency PV. A zero rate
discounts its maturity's whole horizon; it is not an instantaneous forward rate.
At \texttt{risk.py:115--146}, the time is
\begin{equation}
 t_i=\frac{\operatorname{days}(\text{payment}_i-d_s)}{365.0}.
\end{equation}
This counts leap days but always divides by 365, separately from
coupon-based $T_i=q_i/m$. Annual rate times years makes the exponent
dimensionless. The routine uses strictly future payments, adds redemption
at maturity, and returns dirty currency price without subtracting accrual.
It never invokes YTM discounting.

A flat continuous rate numerically equal to $y$ need not give the YTM
price. With common times, conversion $r=m\log(1+y/m)$ would align the
discount factors; the project's different time conventions also preclude
automatic price equality.

\subsection{Independent flat-curve example}
For settlement 1 April 2024, maturity 1 July 2025, $F=100$, $c=0.06$,
$m=2$, and continuous zero rate 0.045:
\begin{center}\small
\begin{tabular}{@{}lrrr@{}}
\toprule
Payment & Calendar days & $t_i$ (years) & $CF_i$\\
\midrule
2024-07-01 & 91 & $91/365$ & 3\\
2025-01-01 & 275 & $275/365$ & 3\\
2025-07-01 & 456 & $456/365$ & 103\\
\bottomrule
\end{tabular}
\end{center}
Explicit summation of
\[
 3e^{-0.045(91/365)}+3e^{-0.045(275/365)}+103e^{-0.045(456/365)}
\]
gave 103.23571061465249. The engine returned the same value, with a
computed difference of 0.0. Dates and cash flows were specified independently.
A constant NSS curve represented the engine input; the reference did
not use the curve-pricing routine.

\subsection{Analytical parallel shock monotonicity}
For a parallel decimal zero-rate shock $s$, fixed times and cash flows give
\begin{equation}
 P(s)=\sum_i CF_i e^{-[r(t_i)+s]t_i},\qquad
 \frac{dP}{ds}=-\sum_i t_iCF_i e^{-[r(t_i)+s]t_i}.
\end{equation}
All terms before the minus sign are nonnegative: $t_i>0$, $CF_i\geq0$,
and exponentials are positive. Redemption ensures one strictly positive
term, even for a zero-coupon bond. Therefore $dP/ds<0$ throughout the
valid real domain. This exact result proves that positive parallel shocks
lower price and negative ones raise it, regardless of the original rate sign.

The risk CLI agrees: base 98.69995920 becomes 94.58065268 under
$+25\bp$ and 103.07256709 under $-25\bp$.
The asymmetric price changes reflect nonlinearity.
Steepener and flattener signs are exposure-dependent and do not follow
from this parallel-shock proof.

\section{Piecewise-linear shocks}
Let $a_1<\cdots<a_L$ be positive anchor maturities and $z_\ell$ their
basis-point shifts. Set decimal shifts $d_\ell=10^{-4}z_\ell$. Then
\begin{equation}
 s(T)=
 \begin{cases}
 d_1,&0<T\leq a_1,\\
 \dfrac{a_{\ell+1}-T}{a_{\ell+1}-a_\ell}d_\ell+
 \dfrac{T-a_\ell}{a_{\ell+1}-a_\ell}d_{\ell+1},
       &a_\ell\leq T\leq a_{\ell+1},\\
 d_L,&T\geq a_L.
 \end{cases}
\end{equation}
For $L=1$ the shock is constant. Substitution at each interval endpoint
gives its anchor shift exactly. Adjacent formulas have matching limits,
proving continuity at every anchor; slopes may change. Tail slopes are
zero. Interpolation weights are dimensionless and sum to one, so $s(T)$
retains decimal rate units.

The \texttt{np.interp} implementation at \texttt{risk.py:160--180}
copies inputs and converts basis points once.
Defaults at lines 183 and 195 are $(-10,0,10)\bp$ for steepening and
$(10,0,-10)\bp$ for flattening, at $(2,10,30)$ years. These are illustrative
scenario definitions, not unique universal market definitions.

Numerical checks also used custom shocks $(-10,0,20)\bp$:
\begin{center}\small
\begin{tabular}{@{}rrrr@{}}
\toprule
$T$ (years) & Custom (bp) & Steepener (bp) & Flattener (bp)\\
\midrule
0.01 & -10 & -10 & 10\\
2 & -10 & -10 & 10\\
6 & -5 & -5 & 5\\
10 & 0 & 0 & 0\\
20 & 10 & 5 & -5\\
30 & 20 & 10 & -10\\
100 & 20 & 10 & -10\\
\bottomrule
\end{tabular}
\end{center}
Anchors, midpoints, and tails agree with independent linear formulas.
For $\eta=10^{-8}$ years, custom one-sided decimal-rate changes were:
\begin{center}\small
\begin{tabular}{@{}rrr@{}}
\toprule
Anchor $a$ & $s(a-\eta)-s(a)$ & $s(a+\eta)-s(a)$\\
\midrule
2 & 0 & $+1.249999886585029\,10^{-12}$\\
10 & $-1.250000103425464\,10^{-12}$ & $+1.000000082740371\,10^{-12}$\\
30 & $-9.999999960041972\,10^{-13}$ & 0\\
\bottomrule
\end{tabular}
\end{center}
These are neighbouring slopes times $\eta$, consistent with continuity.

\section{Key-rate basis and partition of unity}
Take ordered keys $0<k_1<\cdots<k_J$. For $J\geq2$, reconstruct the
endpoint basis functions:
\begin{align}
 b_1(T)&=\begin{cases}
 1,&0<T\leq k_1,\\
 (k_2-T)/(k_2-k_1),&k_1<T<k_2,\\
 0,&T\geq k_2,
 \end{cases}\\
 b_J(T)&=\begin{cases}
 0,&T\leq k_{J-1},\\
 (T-k_{J-1})/(k_J-k_{J-1}),&k_{J-1}<T<k_J,\\
 1,&T\geq k_J.
 \end{cases}
\end{align}
For $1<j<J$,
\begin{equation}
 b_j(T)=\begin{cases}
 (T-k_{j-1})/(k_j-k_{j-1}),&k_{j-1}<T\leq k_j,\\
 (k_{j+1}-T)/(k_{j+1}-k_j),&k_j<T<k_{j+1},\\
 0,&\text{otherwise}.
 \end{cases}
\end{equation}
For a single key, $b_1(T)=1$ everywhere.
The functions are continuous, dimensionless, and lie in $[0,1]$.
Substitution proves $b_j(k_j)=1$ and $b_j(k_\ell)=0$ for $\ell\ne j$.

\subsection{Piecewise proof of partition of unity}
There are four exhaustive kinds of positive maturity.
For $0<T<k_1$, only $b_1$ is nonzero and equals 1.
At $T=k_\ell$, only $b_\ell$ is nonzero and equals 1.
Between adjacent keys $k_\ell<T<k_{\ell+1}$, only the two neighbouring
bases are nonzero, and
\begin{equation}
 b_\ell(T)+b_{\ell+1}(T)
 =\frac{k_{\ell+1}-T}{k_{\ell+1}-k_\ell}
 +\frac{T-k_\ell}{k_{\ell+1}-k_\ell}=1.
\end{equation}
For $T>k_J$, only $b_J$ is nonzero and equals 1.
A single key satisfies the same result by definition.
Consequently the exact identity
\begin{equation}\label{eq:partition}
 \boxed{\sum_{j=1}^{J}b_j(T)=1\quad\text{for every }T>0}
\end{equation}
holds, including both tails.
The source at \texttt{risk.py:218--234} interpolates each identity-matrix
row over keys with constant endpoint tails, reproducing these definitions.

\subsection{Dense independent numerical check}
The public \texttt{key\_rate\_basis} was evaluated for $(2,5,10,30)$
on the sorted unique union of 100,001 linearly spaced maturities from
0.01 to 100, all four keys, and midpoints $(3.5,7.5,20)$.
The union contained 100,008 points. Basic NumPy reductions checked the
entire matrix. A separate piecewise construction using comparisons and
linear formulas, without \texttt{np.interp}, provided a reference.
\begin{center}\small
\begin{tabular}{@{}lr@{}}
\toprule
Check & Actual result\\
\midrule
$\max_T|\sum_j b_j(T)-1|$ & 0.0\\
Minimum / maximum weight & 0.0 / 1.0\\
Maximum difference from independent basis & $1.110223024625157\,10^{-16}$\\
Basis matrix at keys & Exactly the $4\times4$ identity\\
\bottomrule
\end{tabular}
\end{center}
This covers below 2Y, all intervals, and beyond 30Y.
Zero error is an observed binary64 result on this grid; the proof
establishes the identity for all positive maturities.

\section{Key-rate DV01, parallel DV01, and reconciliation}
\subsection{Full repricing and individual checks}
For $h=10^{-4}$, define $r_j^\pm(T)=r(T)\pm h b_j(T)$.
If $P_j^\pm$ are the corresponding fully repriced dirty prices, then
\begin{equation}
 {\rm KRDV01}_j=\frac{P_j^- -P_j^+}{2}.
\end{equation}
Nonnegative weights and positive cash flows give nonnegative down-minus-up
sensitivity; a key with no exposure can give zero.
The code at \texttt{risk.py:237--256} passes plus/minus identity-row
shocks into the basis-point factory, converting to $h$ exactly once.

The independent risk bond used settlement 15 April 2025, maturity
15 January 2055, $F=100$, $c=0.04$, $m=2$.
All 60 dates and cash flows were explicitly constructed.
First and last times were 0.2493150684931507 and
29.77260273972603 years. Base rates were scalar evaluations of the
public fit to the CSV, explicitly interpreted as continuous zeros.
The calibrated curve is the input to this risk check; no risk method
generated reference prices. Independent basis weights and explicit
exponential sums produced all shocked prices. Independent and engine
base prices were both 98.69995919854456.
\begin{center}\small
\begin{tabular}{@{}rrrr@{}}
\toprule
Key (years) & Engine KRDV01 & Independent repricing & Difference\\
\midrule
2 & 0.002369064438994428 & 0.002369064438994428 & 0\\
5 & 0.007200603041120246 & 0.007200603041120246 & 0\\
10 & 0.039682235777171115 & 0.039682235777171115 & 0\\
30 & 0.120474467420841340 & 0.120474467420841340 & 0\\
\bottomrule
\end{tabular}
\end{center}
Units are currency for $F=100$ and a one-basis-point basis move.

\subsection{First-order additivity and finite-bump nonadditivity}
For a decimal key amplitude $u_j$, differentiating at zero shock gives
\[
 -\left.\frac{\partial P}{\partial u_j}\right|_0
 =\sum_i V_i t_i b_j(t_i).
\]
Summation over keys and (\ref{eq:partition}) establish the exact identity
\begin{equation}
 \sum_j\left(-\left.\frac{\partial P}{\partial u_j}\right|_0\right)
 =\sum_i V_i t_i=-\left.\frac{dP}{ds}\right|_0.
\end{equation}
Multiplying by $h$ proves exact additivity of the first-order DV01
approximations.

Full-repricing parallel DV01, at \texttt{risk.py:259--267}, is
\begin{equation}
 {\rm DV01}_{\parallel}=\frac{P(r-h)-P(r+h)}2.
\end{equation}
Here $P(r\pm h)$ prices all payments under the uniformly shifted curve.
Define $\sinh z=(e^z-e^{-z})/2$. Exact real-arithmetic identities are
\begin{align}
 {\rm KRDV01}_j&=\sum_i V_i\sinh(h t_i b_j(t_i)),\\
 {\rm DV01}_{\parallel}&=\sum_i V_i\sinh(h t_i).
\end{align}
For $\Delta=\sum_j{\rm KRDV01}_j-{\rm DV01}_{\parallel}$, expansion
$\sinh z=z+z^3/6+z^5/120+\cdots$ yields
\begin{equation}\label{eq:cubic}
 \Delta=\frac{h^3}{6}\sum_i V_i t_i^3
       \left(\sum_j b_j(t_i)^3-1\right)+O(h^5).
\end{equation}
The linear term cancels by partition of unity. Even powers cancel in
central repricing. The generic leading discrepancy is therefore
$O(h^3)$ in price units, or $O(h^2)$ relative to nonzero parallel DV01
of order $h$.

Since weights lie in $[0,1]$ and sum to one,
$\sum_j b_j^p\leq1$ for every integer $p>1$.
All coefficients in the positive-$z$ series for $\sinh z$ are positive,
so $\Delta\leq0$. It is strictly negative if a positive-PV payment is
split across two keys; it can be zero when every payment has a single
weight equal to one. First-order additivity does not imply finite-bump equality.

Writing $\Delta_3$ for the displayed cubic term,
\begin{equation}
 |\Delta-\Delta_3|\leq
 \sum_i V_i\frac{(h t_i)^5\cosh(h t_i)}{120},
 \qquad \cosh z=(e^z+e^{-z})/2.
\end{equation}
Each omitted odd-power coefficient contains
$1-\sum_j b_j^{2p+1}\in[0,1]$, and the remaining $\sinh$ tail is bounded
by $z^5\cosh z/120$, proving this sufficient bound.
Similarly $0\leq-\Delta\leq\sum_i V_i(h t_i)^3\cosh(h t_i)/6$.
Floating-point repricing adds rounding beyond these real-arithmetic bounds.

\subsection{Actual reconciliation}
\begin{center}\small
\begin{tabular}{@{}lr@{}}
\toprule
Quantity & Actual result\\
\midrule
Parallel DV01 (engine and independent) & 0.16972640002153838\\
Sum of KRDV01s (engine and independent) & 0.16972637067812713\\
$\Delta$ (sum minus parallel) & $-2.9343411256377294\,10^{-8}$\\
Relative $\Delta/{\rm DV01}_{\parallel}$ & $-1.728865471291066\,10^{-7}$\\
First-order parallel sensitivity times $h$ & 0.16972622075771868\\
Derived cubic term $\Delta_3$ & $-2.9343412800122713\,10^{-8}$\\
Independent $\sinh$ difference & $-2.93434242885617\,10^{-8}$\\
Fifth-order remainder bound & $7.086922071550366\,10^{-14}$\\
\bottomrule
\end{tabular}
\end{center}
Cubic prediction, exact-form $\sinh$ computation, and repricing agree
within small truncation and rounding effects.
The relative difference is a dimensionless ratio of sensitivities,
not a basis-point rate change. Additional independent bump checks give
\begin{center}\small
\begin{tabular}{@{}rrr@{}}
\toprule
$h$ & $\Delta$ from $\sinh$ sums & $\Delta/h^3$\\
\midrule
$2\,10^{-4}$ & $-2.3474766955309753\,10^{-7}$ & -29343.458694137185\\
$10^{-4}$ & $-2.9343424288561700\,10^{-8}$ & -29343.424288561695\\
$5\,10^{-5}$ & $-3.6679269569158863\,10^{-9}$ & -29343.415655327084\\
\bottomrule
\end{tabular}
\end{center}
Near constancy of $\Delta/h^3$ supports the cubic leading order.
These extra bumps were temporary checks, not new API parameters.
Curve DV01 and YTM DV01 differentiate different rate variables with
different times and compounding, so their numerical equality is not required.

\section{Dimensional and unit audit}
\small
\begin{longtable}{@{}p{0.25\textwidth}p{0.69\textwidth}@{}}
\toprule
Quantity or coordinate & Audit finding\\
\midrule
\endhead
Currency and price & $F,C,CF_i,PV_i,V_i,AI$, prices, P\&L, and DV01 are amounts for the supplied face value. Default $F=100$ makes the numerical price a price per 100 face; other face values are not automatically normalized.\\
Face normalization & Price per 100 is $100P/F$. Scaling face by $\lambda>0$ scales all cash flows, prices, accrual, and DV01 by $\lambda$, while normalized durations and convexity remain unchanged. This follows exactly from linear sums and their ratios.\\
Decimal rates & $c,y,r,\beta$ use decimals; 0.05 means 5\%. Rate conventions still differ: $y$ is nominal coupon-compounded, zero $r$ is continuous, and generic NSS rates inherit the caller's convention.\\
Percentages & Displayed percentage number is $100r$; 0.045 is displayed as 4.50\%. Plotting explicitly multiplies rates by 100 and labels the percentage axis.\\
Percentage points & Moving from 4\% to 5\% is a change of one percentage point, or 0.01 decimal, or 100 bp. It is a 25\% relative increase in the rate, which is a different concept.\\
Basis points & One bp is $10^{-4}$ decimal; 25 bp is 0.0025. Shock factories multiply bp inputs by $10^{-4}$ once. Returned callables and \texttt{ShockedCurve} use decimal shifts without another conversion.\\
Coupon periods & $q_i$ is a dimensionless fractional period count; $q_i/m$ is coupon-based time in years. $\alpha$ is a ratio of calendar-day counts, not a rate.\\
Curve years & NSS $T,\tau$, shock anchors, and keys are years. Risk times use actual days/365, including leap days in the numerator.\\
Discount factors & $1+y/m$, $(1+y/m)^{-q_i}$, $T/\tau$, and $\exp(-rt)$ have consistent dimensionless bases or arguments.\\
Durations & Macaulay is PV-weighted coupon-based years. Modified duration is $-P'/P$ for annual decimal YTM and equals Macaulay divided by $1+y/m$.\\
Convexity & $P''/P$ with respect to annual decimal YTM, conventionally years squared. No $1/2$ belongs inside this definition.\\
DV01 & Price units for a one-bp move. Central repricing divides the down-minus-up price difference by 2, not by $2h$. Calling it price units per 1 bp specifies the bump, not an unscaled derivative.\\
RMSE & Decimal-rate RMSE is $\sqrt{\sum e_i^2/n}$; bp RMSE multiplies by 10,000, not by 100.\\
Relative reconciliation & $\Delta/{\rm DV01}_{\parallel}$ is dimensionless. Converting that ratio to percent is unrelated to converting a rate to bp.\\
\bottomrule
\end{longtable}
\normalsize
A coordinate change makes hidden scale errors explicit. If $z=100y$ is
the numerical percentage quote, then $dP/dz=P'/100$ and
$d^2P/dz^2=P''/10{,}000$. If $u=10{,}000y$ is the numerical bp quote,
then $dP/du=P'/10{,}000$. The code's derivatives are with respect to
$y$, and the DV01 bump supplies the $10^{-4}$ conversion explicitly.
No hidden factor-of-100 or factor-of-10,000 defect was found.
The conventional year labels for duration and convexity are consistent
with annual-rate perturbations and the factors of $1/m$; they do not
authorize substitution of percentage-number changes into decimal formulas.

\section{Existing tests as evidence and source correspondence}
The complete baseline command, run from the repository root, was
\begin{verbatim}
.\.venv\Scripts\python.exe -m pytest -q
\end{verbatim}
The actual result was:
\begin{verbatim}
234 passed in 4.59s
\end{verbatim}
The process completed successfully. This is a current execution result,
not a README claim.

\small
\begin{longtable}{@{}p{0.20\textwidth}p{0.36\textwidth}p{0.37\textwidth}@{}}
\toprule
Evidence & Inspected coverage & Independent review assessment\\
\midrule
\endhead
\texttt{test\_bonds.py} &
Par, accrual, calendar boundaries, month ends, leap days, fractional periods,
zero coupons, annual/semiannual frequency, negative/zero YTM, face scaling,
scalar precision, invalid inputs, analytics and CLI. &
Existing derivative tests use engine prices at one step.
This report additionally derives derivatives and constructs its own
reference cash flows and price function over seven steps.\\
\texttt{test\_curves.py} &
Formula/shape checks, 70-digit short-tenor references, synthetic fits,
RMSE, deterministic repeats, candidate selection and failure handling,
plots and CLI. &
Existing noiseless fits usually generate observations with the model API.
This review instead generates them from 90-digit independent exponentials
and also compares a dense between-observation grid.\\
\texttt{test\_risk.py} &
Explicit flat-curve prices, leap-day timing, settlement exclusion,
shock signs and units, interpolation/tails, custom/single keys,
reconciliation bounds, scaling, immutability, validation, CLI. &
The tests contain substantive financial checks, not placeholders.
This review independently proves all-maturity partition of unity,
checks 100,008 points, and reprices all individual key sensitivities.\\
\bottomrule
\end{longtable}
\normalsize

Source locations refer to the audited revision:
\small
\begin{longtable}{@{}p{0.36\textwidth}p{0.58\textwidth}@{}}
\toprule
Location & Mathematical correspondence\\
\midrule
\endhead
\texttt{bonds.py:76--136} & Validated bond domain, maturity-anchored schedule, coupon-date exclusion, period accrual, positive discount base.\\
\texttt{bonds.py:139--217} & Fractional-period PV sum, clean/dirty identity, analytical durations and convexity, central YTM DV01.\\
\texttt{curves.py:29--93} & NSS loading series, stable evaluation, parameter and maturity domain.\\
\texttt{curves.py:120--214} & RMSE conversion, observations, deterministic constrained multi-start fit, candidate selection.\\
\texttt{curves.py:217--246} & Plot converts decimals to labelled percentages; no financial rate conversion.\\
\texttt{risk.py:71--146} & Additive decimal curve shifts, scenario P\&L, Actual/365 continuous discounting.\\
\texttt{risk.py:149--215} & Single bp conversion, linear shocks and constant tails, default scenarios, full repricing.\\
\texttt{risk.py:218--267} & Key basis, key-rate and parallel central DV01.\\
\texttt{cli.py:18--135} & Invented CSV fit, explicit zero interpretation for risk, illustrative commands and displayed units.\\
\texttt{\_\_init\_\_.py} & Exports existing bond, NSS, and risk APIs without alternate financial formulas.\\
\texttt{README.md} & Documented formulas, units, and limitations agree with the reviewed mathematics.\\
\texttt{pyproject.toml} & Python 3.12+; declared dependencies and argparse-based CLI entry point. No configuration changes made.\\
\bottomrule
\end{longtable}
\normalsize
The existing tests and the additional checks agree with the derivations.
Passing tests remain finite evidence, not proofs of universal numerical
behaviour.

\section{Actual CLI execution record}
The following commands were executed from the repository root after the
baseline suite:
\begin{verbatim}
.\.venv\Scripts\python.exe -m fixed_income.cli bond
.\.venv\Scripts\python.exe -m fixed_income.cli curve
.\.venv\Scripts\python.exe -m fixed_income.cli risk
.\.venv\Scripts\python.exe -m fixed_income.cli demo
\end{verbatim}
All completed. The tables retain their actual displayed numerical
precision. The \texttt{demo} command reproduced the bond, curve, and
risk outputs below in that order, with the same numerical results;
its curve fit is reused for its risk section.

\subsection{Bond command}
The printed example was explicitly invented. Settlement was 2025-04-15,
maturity 2030-01-15, face 100.00 currency units, coupon 4.00\%
(0.040000 decimal), nominal annual YTM 4.50\% (0.045000 decimal),
frequency 2 payments per year.
\begin{center}
\begin{tabular}{@{}lr@{}}
\toprule
Printed measure & Printed value\\
\midrule
Accrued interest (currency) & 0.994475\\
Dirty price (currency) & 98.871313\\
Clean price (currency) & 97.876838\\
Macaulay duration (years) & 4.326727\\
Modified duration (years) & 4.231518\\
Convexity (years squared) & 21.110445\\
DV01 (currency for 1 bp) & 0.041838\\
\bottomrule
\end{tabular}
\end{center}
The CLI explicitly labels central dirty-price repricing and
$1\bp=0.0001$ decimal YTM.

\subsection{Curve command and inspected CSV}
The actual CSV observations were:
\begin{center}\small
\begin{tabular}{@{}rrrrr@{}}
\toprule
Maturity (years) & Rate (decimal) && Maturity (years) & Rate (decimal)\\
\midrule
0.25 & 0.0420 && 5 & 0.0370\\
0.5 & 0.0410 && 7 & 0.0380\\
1 & 0.0390 && 10 & 0.0390\\
2 & 0.0370 && 20 & 0.0410\\
3 & 0.0365 && 30 & 0.0415\\
\bottomrule
\end{tabular}
\end{center}
These are invented observations, not market data. Their fit does not
bootstrap a zero curve. The command printed:
\begin{center}
\begin{tabular}{@{}lr@{}}
\toprule
Measure & Printed value\\
\midrule
$\beta_0$ (decimal rate) & 0.0395973473\\
$\beta_1$ (decimal rate) & 0.0040652482\\
$\beta_2$ (decimal rate) & -0.0190419619\\
$\beta_3$ (decimal rate) & 0.0091397288\\
$\tau_1$ (years) & 1.6287979110\\
$\tau_2$ (years) & 16.1937672075\\
Optimizer success & True\\
Optimizer status & 2\\
Selected-start function evaluations & 17\\
Successful starts & 6/6\\
RMSE (decimal rate) & 0.0000782762\\
RMSE (basis points) & 0.782762\\
\bottomrule
\end{tabular}
\end{center}
The termination message stated that the \texttt{ftol} termination
condition was satisfied. The plot path printed was
\texttt{outputs\textbackslash yield\_curve.png}.
Curve and demo wrote the existing gitignored plot location.
Rounded display values need not retain exact arithmetic conversion
between their last printed digits; the underlying RMSE uses one
factor of 10,000.

\subsection{Risk command}
The printed risk example used settlement 2025-04-15, maturity 2055-01-15,
face 100.00, coupon 4.00\% annually, and semiannual payments.
It explicitly interpreted fitted rates as continuous zeros for this
educational example, with Actual/365 times.
All price, P\&L, and DV01 values below are currency for that face value:
\begin{center}\small
\begin{tabular}{@{}lrrr@{}}
\toprule
Scenario & Base & Shocked & P\&L\\
\midrule
Parallel $+25\bp$ & 98.69995920 & 94.58065268 & -4.11930652\\
Parallel $-25\bp$ & 98.69995920 & 103.07256709 & +4.37260789\\
Steepener $-10/0/+10\bp$ & 98.69995920 & 97.57938095 & -1.12057825\\
Flattener $+10/0/-10\bp$ & 98.69995920 & 99.85176589 & +1.15180669\\
\bottomrule
\end{tabular}
\end{center}
Shape anchors were 2/10/30 years. Scenario P\&L is exactly defined as
shocked minus base price; subtracting the rounded displayed prices can
differ in the last digit from the displayed full-precision calculation.
\begin{center}\small
\begin{tabular}{@{}lr@{}}
\toprule
Sensitivity & Printed value\\
\midrule
2Y KRDV01 & 0.0023690644\\
5Y KRDV01 & 0.0072006030\\
10Y KRDV01 & 0.0396822358\\
30Y KRDV01 & 0.1204744674\\
Parallel DV01 & 0.1697264000\\
Sum of KRDV01 & 0.1697263707\\
Reconciliation (sum minus parallel) & $-2.934341125638\,10^{-8}$\\
\bottomrule
\end{tabular}
\end{center}
The printed DV01 definition was $[P(-1\bp)-P(+1\bp)]/2$, with one bp
equal to 0.0001 decimal zero rate.

\section{Verification conclusions and document build}
No mathematical defect was found in the reviewed bond pricing, accrual,
duration, convexity, DV01, NSS loading/calibration definitions,
continuous-zero pricing, scenario interpolation, or key-rate basis and
reconciliation formulas. No unresolved discrepancy requires a source
correction. Exact identities have been derived under explicit assumptions;
sample numerical agreement is presented as supporting evidence.

The principal approximation effects are central derivative truncation
of order $h^2$, finite-bump DV01 corrections of order $h^3$,
NSS small-argument polynomial remainders of order $x^6$, floating-point
rounding, and numerical optimization. First-order key-rate additivity
coexists with a generally nonzero cubic full-repricing discrepancy.
No missing factor of 100 or 10,000 was identified.

The limitations are those of the stated conventions and finite numerical
evidence: regular coupon schedules, prescribed accrual and time bases,
caller-assigned rate meaning, no calibration uniqueness/global-minimum
guarantee, and no verification of every extreme floating-point input.
These limitations do not contradict the validated mathematical identities.

\BuildResult
The PowerShell build commands are:
\begingroup\small
\begin{verbatim}
latexmk --version
Push-Location .\paper
latexmk -pdf -interaction=nonstopmode -halt-on-error verification_report.tex
Pop-Location
\end{verbatim}
\endgroup
The source report is \texttt{paper/verification\_report.tex}.
Only this source and a minimal paper-specific ignore file are intended
for the documentation commit; PDF and LaTeX auxiliary outputs are excluded.

\appendix
\section{Independent stdin-fed numerical scripts}\label{sec:scripts}
These are the executed review scripts. Run each from the repository root
inside the following PowerShell wrapper:
\begin{verbatim}
@'
# Paste one Python script below here.
'@ | .\.venv\Scripts\python.exe -
\end{verbatim}
The scripts import only standard-library modules, NumPy, and public
\texttt{fixed\_income} APIs. Reference calculations do not call the
analytics method being checked. Display-only decimal truncation in printed
tables is separate from the full binary64 values used in calculations.

\subsection{Bond cash flows, derivatives, DV01, and coupon boundaries}
\begingroup\scriptsize
\begin{verbatim}
from datetime import date
from math import fsum
import numpy as np
from fixed_income import FixedRateBond
np.set_printoptions(precision=16)
s = date(2025,4,15)
dates = [date(2025,7,15)] + [
    date(yr,mo,15) for yr in range(2026,2030) for mo in (1,7)
] + [date(2030,1,15)]
alpha = (s-date(2025,1,15)).days / (
    date(2025,7,15)-date(2025,1,15)).days
q = 1-alpha+np.arange(len(dates))
cf = np.full(len(dates),2.0)
cf[-1] += 100
def price(y):
    return fsum(float(a)*(1+float(y)/2)**(-float(b))
                for a,b in zip(cf,q))
y = 0.045
p = price(y)
pv = cf*(1+y/2)**(-q)
d1 = -fsum(float(v*t/(2*(1+y/2))) for v,t in zip(pv,q))
d2 = fsum(float(v*t*(t+1)/(4*(1+y/2)**2))
           for v,t in zip(pv,q))
dm = -d1/p
cv = d2/p
b = FixedRateBond(s,dates[-1],0.04)
assert b.coupon_schedule() == tuple(dates)
print("alpha,w,AI",alpha,1-alpha,2*alpha)
print("count,q_first,q_last",len(dates),q[0],q[-1])
print("independent dirty,clean",p,p-2*alpha)
print("engine dirty,clean",b.dirty_price(y),b.clean_price(y))
print("independent Mac,Mod,dP,d2P,conv",dm*(1+y/2),dm,d1,d2,cv)
print("engine Mac,Mod,conv",b.macaulay_duration(y),b.modified_duration(y),b.convexity(y))
print("h,fd_mod,mod_error,fd_second,fd_conv,conv_error")
for h in (1e-2,1e-3,1e-4,1e-5,1e-6,1e-7,1e-8):
    pp,pm = price(y+h),price(y-h)
    fd = -(pp-pm)/(2*h*p)
    sec = (pp-2*p+pm)/h**2
    print(f"{h:.0e},{fd:.15g},{fd-dm:+.9e},{sec:.15g},{sec/p:.15g},{sec/p-cv:+.9e}")
h=1e-4
direct=(price(y-h)-price(y+h))/2
approx=p*dm*h
print("DV01 engine,direct,approx,direct_minus_approx",b.dv01(y),direct,approx,direct-approx)
print("Pminus,Pplus",price(y-h),price(y+h))
print("par m,independent,engine")
for m in (1,2):
    n=5*m
    ref=fsum((100*0.04/m+(100 if i==n else 0)) /
             (1+0.04/m)**i for i in range(1,n+1))
    obj=FixedRateBond(date(2025,1,15),date(2030,1,15),0.04,frequency=m)
    print(m,ref,obj.dirty_price(0.04))
print("boundary settlement,previous,next,n,AI,q_first")
for day in (14,15,16):
    obj=FixedRateBond(date(2025,7,day),date(2027,1,15),0.04)
    a=(obj.settlement_date-obj.previous_coupon_date).days/(
        obj.next_coupon_date-obj.previous_coupon_date).days
    print(obj.settlement_date,obj.previous_coupon_date,obj.next_coupon_date,
          len(obj.coupon_schedule()),obj.accrued_interest(),1-a)
\end{verbatim}
\endgroup
\subsection{NSS high-precision loadings and independent synthetic calibration}
\begingroup\scriptsize
\begin{verbatim}
from decimal import Decimal, localcontext
from math import exp
import numpy as np
from fixed_income import NSSCurve,fit_nss
theta=(0.04,-0.025,0.03,-0.012,1.4,6.5)
def reference(t,parameters=theta):
    with localcontext() as ctx:
        ctx.prec=90
        T=Decimal(str(float(t)))
        b0,b1,b2,b3,t1,t2=map(lambda x:Decimal(str(x)),parameters)
        x1,x2=T/t1,T/t2
        e1,e2=(-x1).exp(),(-x2).exp()
        a1,a2=(1-e1)/x1,(1-e2)/x2
        return float(b0+b1*a1+b2*(a1-e1)+b3*(a2-e2))
known=NSSCurve(*theta)
print("T,engine_rate,reference_error,isolated_L2,relative_L2_error,naive_L1")
for t in (1e-24,1e-18,1e-12,1e-8,0.000999,0.001,0.001001):
    actual=known.yield_rate(t)
    p=(0,0,1,0,1,2)
    iso=NSSCurve(*p).yield_rate(t)
    ref=reference(t,p)
    print(f"{t:.9g},{actual:.17g},{actual-reference(t):+.3e},"
          f"{iso:.17g},{(iso-ref)/ref:+.3e},{(1-exp(-t))/t:.17g}")
print("long T,rate,rate_minus_beta0")
for t in (100,1e4,1e8,1e12):
    r=known.yield_rate(t)
    print(t,repr(r),repr(r-theta[0]))
t=np.geomspace(0.05,40.0,40)
rates=np.array([reference(v) for v in t])
fit=fit_nss(t,rates)
res=rates-fit.curve.yield_rate(t)
grid=np.geomspace(0.05,40.0,1001)
truth=np.array([reference(v) for v in grid])
error=fit.curve.yield_rate(grid)-truth
print("known parameters",theta)
print("fitted",fit.curve)
print("max_abs_observed,rmse_independent,rmse_API,rmse_bps",
      max(abs(res)),np.sqrt(np.mean(res**2)),fit.rmse,fit.rmse_bps)
print("grid max_abs",max(abs(error)))
print("success,status,nfev,starts",fit.success,fit.status,fit.nfev,
      fit.starts_succeeded,fit.starts_attempted)
print("residual_API_max_error",max(abs(fit.residuals-res)))
print("isolated loading maximum relative errors")
tiny=np.array([1e-24,1e-18,1e-12,1e-8,0.000999,
               0.001,0.001001,0.001998,0.002,0.002002])
for index in (1,2,3):
    p=[0,0,0,0,1,2]
    p[index]=1
    ref=np.array([reference(v,p) for v in tiny])
    act=NSSCurve(*p).yield_rate(tiny)
    print(index,float(np.max(np.abs((act-ref)/ref))))
\end{verbatim}
\endgroup
\subsection{Zero prices, shocks, basis, and full-repricing reconciliation}
\begingroup\scriptsize
\begin{verbatim}
from datetime import date
from math import exp,fsum,sinh,cosh
from pathlib import Path
import numpy as np
from fixed_income import (
    FixedRateBond,NSSCurve,fit_nss,curve_price,key_rate_basis,
    key_rate_dv01,parallel_dv01,piecewise_linear_shock,
    steepener_shock,flattener_shock,
)
s=date(2024,4,1)
dates=[date(2024,7,1),date(2025,1,1),date(2025,7,1)]
cf=np.array([3.0,3.0,103.0])
t=np.array([(v-s).days/365 for v in dates])
flat=fsum(float(a)*exp(-0.045*float(v)) for a,v in zip(cf,t))
obj=FixedRateBond(s,dates[-1],0.06)
actual=curve_price(obj,NSSCurve(0.045,0,0,0,1,5))
print("flat days,t,CF",(np.array(t)*365).round().astype(int),t,cf)
print("flat price reference,engine,difference",flat,actual,actual-flat)
keys=np.array([2.0,5.0,10.0,30.0])
def basis(t):
    t=np.asarray(t)
    out=np.zeros(t.shape+(len(keys),))
    out[...,0]=np.where(t<=keys[0],1,0)
    out[...,-1]=np.where(t>=keys[-1],1,0)
    for j,(left,right) in enumerate(zip(keys[:-1],keys[1:])):
        mask=(t>left)&(t<right)
        out[...,j]+=np.where(mask,(right-t)/(right-left),0)
        out[...,j+1]+=np.where(mask,(t-left)/(right-left),0)
    for j,k in enumerate(keys[1:-1],1):
        out[...,j]+=np.where(t==k,1,0)
    return out
grid=np.unique(np.r_[np.linspace(0.01,100,100001),keys,[3.5,7.5,20]])
weights=key_rate_basis(grid)
print("basis n,max_partition,min,max,max_reference_error",
      grid.size,np.max(np.abs(weights.sum(axis=1)-1)),
      weights.min(),weights.max(),np.max(np.abs(weights-basis(grid))))
print("basis_at_keys",key_rate_basis(keys).tolist())
anchors=[2.,10.,30.]
shifts=[-10.,0.,20.]
shock=piecewise_linear_shock(anchors,shifts)
print("T,generic_bp,steepener_bp,flattener_bp")
for v in (0.01,2,6,10,20,30,100):
    print(v,shock(v)*1e4,steepener_shock()(v)*1e4,
          flattener_shock()(v)*1e4)
eps=1e-8
print("anchor,left_minus_anchor,right_minus_anchor (decimal)")
for v in anchors:
    print(v,shock(v-eps)-shock(v),shock(v+eps)-shock(v))
source=Path("examples")/"illustrative_curve.csv"
data=np.genfromtxt(source,delimiter=",",names=True)
curve=fit_nss(data["maturity_years"],data["rate"]).curve
s=date(2025,4,15)
dates=[date(2025,7,15)]+[
    date(yr,mo,15) for yr in range(2026,2055) for mo in (1,7)
]+[date(2055,1,15)]
cf=np.full(len(dates),2.0)
cf[-1]+=100
t=np.array([(v-s).days/365 for v in dates])
rates=np.array([curve.yield_rate(float(v)) for v in t])
w=basis(t)
pv=np.array([a*exp(-r*v) for a,r,v in zip(cf,rates,t)])
def price(shifts):
    return fsum(float(a)*exp(-float(r+dv)*float(v))
                for a,r,dv,v in zip(cf,rates,np.broadcast_to(shifts,t.shape),t))
bond=FixedRateBond(s,dates[-1],0.04)
assert bond.coupon_schedule()==tuple(dates)
print("risk count,first_time,last_time,independent_base,engine_base",
      len(dates),t[0],t[-1],price(0),curve_price(bond,curve))
h=1e-4
ind=np.array([(price(-h*w[:,j])-price(h*w[:,j]))/2
              for j in range(len(keys))])
implemented=key_rate_dv01(bond,curve)
print("key,engine,independent,difference")
for item,v in zip(implemented,ind):
    print(item.tenor_years,repr(item.dv01),repr(v),repr(item.dv01-v))
parallel=parallel_dv01(bond,curve)
total=fsum(item.dv01 for item in implemented)
pind=(price(-h)-price(h))/2
cubic=fsum(float(p)*(h*float(v))**3*(fsum(float(a)**3 for a in row)-1)/6
           for p,v,row in zip(pv,t,w))
bound=fsum(float(p)*(h*float(v))**5*cosh(h*float(v))/120
           for p,v in zip(pv,t))
stable=fsum(float(p)*(fsum(sinh(h*float(v)*float(a)) for a in row)
                        -sinh(h*float(v))) for p,v,row in zip(pv,t,w))
print("parallel_API,parallel_independent,sum_API,sum_independent",
      repr(parallel),repr(pind),repr(total),repr(fsum(ind)))
print("difference,relative,cubic,stable,fifth_bound",
      repr(total-parallel),repr((total-parallel)/parallel),
      repr(cubic),repr(stable),repr(bound))
print("first_order_parallel",fsum(float(p*v*h) for p,v in zip(pv,t)))
print("h,stable_difference,difference_over_hcubed")
for h in (2e-4,1e-4,5e-5):
    d=fsum(float(p)*(fsum(sinh(h*float(v)*float(a)) for a in row)
                    -sinh(h*float(v))) for p,v,row in zip(pv,t,w))
    print(h,repr(d),repr(d/h**3))
\end{verbatim}
\endgroup
\end{document}
